\providecommand{\Res}[2]{\qty[round-mode=figures,round-precision=3]{\fpeval{#1}}{#2}}
\providecommand{\Val}[1]{\num[round-mode=figures,round-precision=4]{\fpeval{#1}}}
\providecommand{\ResSci}[2]{\qty[round-mode=figures,round-precision=3,exponent-mode=scientific]{\fpeval{#1}}{#2}}
\subsubsection*{Question a) Fréquence perçue}
\[ f_{\text{moy}} = \frac{f_1 + f_2}{2} = \frac{\num{\v{f1}} + \num{\fpeval{\v{f1}+\v{df}}}}{2} = \Res{\v{f1} + \v{df}/2}{Hz} \]
On perçoit un son de \textbf{\Res{\v{f1} + \v{df}/2}{Hz}}.

\subsubsection*{Question b) Période des battements}
L'amplitude $2A\cos(2\pi\,\tfrac{\Delta f}{2}\,t)$ oscille à $\Delta f/2$, mais l'intensité dépend du carré de l'amplitude : elle est maximale à la fois quand le cosinus vaut $+1$ et $-1$, donc deux fois par période de l'enveloppe. La fréquence des battements est donc $f_{\text{batt}} = f_2 - f_1 = \qty{\v{df}}{Hz}$.
\[ T_{\text{batt}} = \frac{1}{f_{\text{batt}}} = \frac{1}{\num{\v{df}}} = \Res{1/\v{df}}{s} \]
Deux maximums consécutifs sont séparés de \textbf{\Res{1/\v{df}}{s}}.

\subsubsection*{Question c) Nombre de battements}
\[ N = f_{\text{batt}}\,\Delta t = \num{\v{df}} \cdot \num{\v{t}} = \fpeval{\v{df}*\v{t}} \]
On entend \textbf{\fpeval{\v{df}*\v{t}} battements}.
